%% How a vision-language model reads a page: patch embedding -> vision encoder ->
%% projection -> autoregressive decoder conditioned on visual tokens + prompt.
%%
%% Hand-authored TikZ (not generated): it depicts the standard architecture described
%% in Section~\ref{sec:bg-vlm} (Vaswani et al. 2017; Dosovitskiy et al. 2021), not a
%% measurement, so there is no data source to derive it from. It deliberately carries
%% no invoice-specific content: nothing in this architecture knows what an invoice is,
%% which is exactly the point the accompanying prose makes.
%%
%% Input from a chapter with:  %% How a vision-language model reads a page: patch embedding -> vision encoder ->
%% projection -> autoregressive decoder conditioned on visual tokens + prompt.
%%
%% Hand-authored TikZ (not generated): it depicts the standard architecture described
%% in Section~\ref{sec:bg-vlm} (Vaswani et al. 2017; Dosovitskiy et al. 2021), not a
%% measurement, so there is no data source to derive it from. It deliberately carries
%% no invoice-specific content: nothing in this architecture knows what an invoice is,
%% which is exactly the point the accompanying prose makes.
%%
%% Input from a chapter with:  %% How a vision-language model reads a page: patch embedding -> vision encoder ->
%% projection -> autoregressive decoder conditioned on visual tokens + prompt.
%%
%% Hand-authored TikZ (not generated): it depicts the standard architecture described
%% in Section~\ref{sec:bg-vlm} (Vaswani et al. 2017; Dosovitskiy et al. 2021), not a
%% measurement, so there is no data source to derive it from. It deliberately carries
%% no invoice-specific content: nothing in this architecture knows what an invoice is,
%% which is exactly the point the accompanying prose makes.
%%
%% Input from a chapter with:  %% How a vision-language model reads a page: patch embedding -> vision encoder ->
%% projection -> autoregressive decoder conditioned on visual tokens + prompt.
%%
%% Hand-authored TikZ (not generated): it depicts the standard architecture described
%% in Section~\ref{sec:bg-vlm} (Vaswani et al. 2017; Dosovitskiy et al. 2021), not a
%% measurement, so there is no data source to derive it from. It deliberately carries
%% no invoice-specific content: nothing in this architecture knows what an invoice is,
%% which is exactly the point the accompanying prose makes.
%%
%% Input from a chapter with:  \input{figures/vlm-anatomy.tex}

\begin{tikzpicture}[
    font=\small,
    node distance=6mm,
    block/.style={
      draw=horusaccent, line width=0.8pt, rounded corners=2pt,
      fill=horuswash, text width=2.3cm, align=center,
      minimum height=1.5cm, inner sep=3pt,
    },
    endpoint/.style={
      draw=horusmuted, line width=0.6pt, rounded corners=1pt,
      fill=white, text width=1.4cm, align=center,
      font=\scriptsize, minimum height=1.0cm, inner sep=3pt,
    },
    flow/.style={-{Stealth[length=2.4mm]}, line width=0.8pt, draw=horusaccent},
    onflow/.style={font=\scriptsize, text=horusink, align=center, inner sep=2pt},
  ]

  %% Page -> patch grid. The grid is drawn, not described: the one visual idea a
  %% reader must take away is that the page becomes a sequence of image tokens.
  %% Labelled with the concrete input ("invoice page", not the abstract "page
  %% image") per the examiner's 2026-08-18 note, registry row R08.
  \node[endpoint, minimum height=1.5cm] (page) {invoice\\page};
  \begin{scope}[shift={($(page.east)+(7mm,-5.4mm)$)}]
    \foreach \x in {0,1,2} \foreach \y in {0,1,2}
      \draw[draw=horusmuted, line width=0.5pt, fill=white]
        (\x*3.4mm, \y*3.4mm) rectangle ++(3.0mm,3.0mm);
  \end{scope}
  \node[onflow, anchor=north] at ($(page.east)+(12mm,-7mm)$) {fixed-size\\patches};

  \node[block, right=22mm of page] (vit) {\textbf{Vision encoder}\\[2pt]\scriptsize a
    transformer over the patch sequence};
  \node[block, right=10mm of vit] (proj) {\textbf{Projection}\\[2pt]\scriptsize maps visual
    embeddings into the language model's space};
  \node[block, right=10mm of proj] (dec) {\textbf{Decoder}\\[2pt]\scriptsize
    autoregressively generates text conditioned on visual tokens and the prompt};
  \node[endpoint, right=4mm of dec] (out) {output\\tokens};

  \draw[flow] (page) -- node[onflow, above=1mm, pos=0.5] {} (vit);
  \draw[flow] (vit) -- (proj);
  \draw[flow] (proj) -- (dec);
  \draw[flow] (dec) -- (out);

  \node[onflow, anchor=south] at ($(vit.north east)!0.5!(proj.north west) + (0,2.5mm)$)
    {visual\\embeddings};
  \node[onflow, anchor=south] at ($(proj.north east)!0.5!(dec.north west) + (0,2.5mm)$)
    {visual\\tokens};

  %% The prompt enters the decoder from below: it is the second conditioning channel,
  %% and everything task-specific in this thesis lives in it.
  \node[endpoint, below=9mm of dec] (prompt) {textual\\prompt};
  \draw[flow] (prompt) -- (dec);

  \node[
    draw=horusmuted, line width=0.5pt, rounded corners=2pt, dotted,
    fill=white, text=horusink, font=\scriptsize, align=left,
    text width=4.6cm, inner sep=4pt,
    below=9mm of vit.south west, anchor=north west,
  ] {No component of this architecture is invoice-specific: all task behaviour
    is elicited by the prompt and learned from training data.};

\end{tikzpicture}


\begin{tikzpicture}[
    font=\small,
    node distance=6mm,
    block/.style={
      draw=horusaccent, line width=0.8pt, rounded corners=2pt,
      fill=horuswash, text width=2.3cm, align=center,
      minimum height=1.5cm, inner sep=3pt,
    },
    endpoint/.style={
      draw=horusmuted, line width=0.6pt, rounded corners=1pt,
      fill=white, text width=1.4cm, align=center,
      font=\scriptsize, minimum height=1.0cm, inner sep=3pt,
    },
    flow/.style={-{Stealth[length=2.4mm]}, line width=0.8pt, draw=horusaccent},
    onflow/.style={font=\scriptsize, text=horusink, align=center, inner sep=2pt},
  ]

  %% Page -> patch grid. The grid is drawn, not described: the one visual idea a
  %% reader must take away is that the page becomes a sequence of image tokens.
  %% Labelled with the concrete input ("invoice page", not the abstract "page
  %% image") per the examiner's 2026-08-18 note, registry row R08.
  \node[endpoint, minimum height=1.5cm] (page) {invoice\\page};
  \begin{scope}[shift={($(page.east)+(7mm,-5.4mm)$)}]
    \foreach \x in {0,1,2} \foreach \y in {0,1,2}
      \draw[draw=horusmuted, line width=0.5pt, fill=white]
        (\x*3.4mm, \y*3.4mm) rectangle ++(3.0mm,3.0mm);
  \end{scope}
  \node[onflow, anchor=north] at ($(page.east)+(12mm,-7mm)$) {fixed-size\\patches};

  \node[block, right=22mm of page] (vit) {\textbf{Vision encoder}\\[2pt]\scriptsize a
    transformer over the patch sequence};
  \node[block, right=10mm of vit] (proj) {\textbf{Projection}\\[2pt]\scriptsize maps visual
    embeddings into the language model's space};
  \node[block, right=10mm of proj] (dec) {\textbf{Decoder}\\[2pt]\scriptsize
    autoregressively generates text conditioned on visual tokens and the prompt};
  \node[endpoint, right=4mm of dec] (out) {output\\tokens};

  \draw[flow] (page) -- node[onflow, above=1mm, pos=0.5] {} (vit);
  \draw[flow] (vit) -- (proj);
  \draw[flow] (proj) -- (dec);
  \draw[flow] (dec) -- (out);

  \node[onflow, anchor=south] at ($(vit.north east)!0.5!(proj.north west) + (0,2.5mm)$)
    {visual\\embeddings};
  \node[onflow, anchor=south] at ($(proj.north east)!0.5!(dec.north west) + (0,2.5mm)$)
    {visual\\tokens};

  %% The prompt enters the decoder from below: it is the second conditioning channel,
  %% and everything task-specific in this thesis lives in it.
  \node[endpoint, below=9mm of dec] (prompt) {textual\\prompt};
  \draw[flow] (prompt) -- (dec);

  \node[
    draw=horusmuted, line width=0.5pt, rounded corners=2pt, dotted,
    fill=white, text=horusink, font=\scriptsize, align=left,
    text width=4.6cm, inner sep=4pt,
    below=9mm of vit.south west, anchor=north west,
  ] {No component of this architecture is invoice-specific: all task behaviour
    is elicited by the prompt and learned from training data.};

\end{tikzpicture}


\begin{tikzpicture}[
    font=\small,
    node distance=6mm,
    block/.style={
      draw=horusaccent, line width=0.8pt, rounded corners=2pt,
      fill=horuswash, text width=2.3cm, align=center,
      minimum height=1.5cm, inner sep=3pt,
    },
    endpoint/.style={
      draw=horusmuted, line width=0.6pt, rounded corners=1pt,
      fill=white, text width=1.4cm, align=center,
      font=\scriptsize, minimum height=1.0cm, inner sep=3pt,
    },
    flow/.style={-{Stealth[length=2.4mm]}, line width=0.8pt, draw=horusaccent},
    onflow/.style={font=\scriptsize, text=horusink, align=center, inner sep=2pt},
  ]

  %% Page -> patch grid. The grid is drawn, not described: the one visual idea a
  %% reader must take away is that the page becomes a sequence of image tokens.
  %% Labelled with the concrete input ("invoice page", not the abstract "page
  %% image") per the examiner's 2026-08-18 note, registry row R08.
  \node[endpoint, minimum height=1.5cm] (page) {invoice\\page};
  \begin{scope}[shift={($(page.east)+(7mm,-5.4mm)$)}]
    \foreach \x in {0,1,2} \foreach \y in {0,1,2}
      \draw[draw=horusmuted, line width=0.5pt, fill=white]
        (\x*3.4mm, \y*3.4mm) rectangle ++(3.0mm,3.0mm);
  \end{scope}
  \node[onflow, anchor=north] at ($(page.east)+(12mm,-7mm)$) {fixed-size\\patches};

  \node[block, right=22mm of page] (vit) {\textbf{Vision encoder}\\[2pt]\scriptsize a
    transformer over the patch sequence};
  \node[block, right=10mm of vit] (proj) {\textbf{Projection}\\[2pt]\scriptsize maps visual
    embeddings into the language model's space};
  \node[block, right=10mm of proj] (dec) {\textbf{Decoder}\\[2pt]\scriptsize
    autoregressively generates text conditioned on visual tokens and the prompt};
  \node[endpoint, right=4mm of dec] (out) {output\\tokens};

  \draw[flow] (page) -- node[onflow, above=1mm, pos=0.5] {} (vit);
  \draw[flow] (vit) -- (proj);
  \draw[flow] (proj) -- (dec);
  \draw[flow] (dec) -- (out);

  \node[onflow, anchor=south] at ($(vit.north east)!0.5!(proj.north west) + (0,2.5mm)$)
    {visual\\embeddings};
  \node[onflow, anchor=south] at ($(proj.north east)!0.5!(dec.north west) + (0,2.5mm)$)
    {visual\\tokens};

  %% The prompt enters the decoder from below: it is the second conditioning channel,
  %% and everything task-specific in this thesis lives in it.
  \node[endpoint, below=9mm of dec] (prompt) {textual\\prompt};
  \draw[flow] (prompt) -- (dec);

  \node[
    draw=horusmuted, line width=0.5pt, rounded corners=2pt, dotted,
    fill=white, text=horusink, font=\scriptsize, align=left,
    text width=4.6cm, inner sep=4pt,
    below=9mm of vit.south west, anchor=north west,
  ] {No component of this architecture is invoice-specific: all task behaviour
    is elicited by the prompt and learned from training data.};

\end{tikzpicture}


\begin{tikzpicture}[
    font=\small,
    node distance=6mm,
    block/.style={
      draw=horusaccent, line width=0.8pt, rounded corners=2pt,
      fill=horuswash, text width=2.3cm, align=center,
      minimum height=1.5cm, inner sep=3pt,
    },
    endpoint/.style={
      draw=horusmuted, line width=0.6pt, rounded corners=1pt,
      fill=white, text width=1.4cm, align=center,
      font=\scriptsize, minimum height=1.0cm, inner sep=3pt,
    },
    flow/.style={-{Stealth[length=2.4mm]}, line width=0.8pt, draw=horusaccent},
    onflow/.style={font=\scriptsize, text=horusink, align=center, inner sep=2pt},
  ]

  %% Page -> patch grid. The grid is drawn, not described: the one visual idea a
  %% reader must take away is that the page becomes a sequence of image tokens.
  %% Labelled with the concrete input ("invoice page", not the abstract "page
  %% image") per the examiner's 2026-08-18 note, registry row R08.
  \node[endpoint, minimum height=1.5cm] (page) {invoice\\page};
  \begin{scope}[shift={($(page.east)+(7mm,-5.4mm)$)}]
    \foreach \x in {0,1,2} \foreach \y in {0,1,2}
      \draw[draw=horusmuted, line width=0.5pt, fill=white]
        (\x*3.4mm, \y*3.4mm) rectangle ++(3.0mm,3.0mm);
  \end{scope}
  \node[onflow, anchor=north] at ($(page.east)+(12mm,-7mm)$) {fixed-size\\patches};

  \node[block, right=22mm of page] (vit) {\textbf{Vision encoder}\\[2pt]\scriptsize a
    transformer over the patch sequence};
  \node[block, right=10mm of vit] (proj) {\textbf{Projection}\\[2pt]\scriptsize maps visual
    embeddings into the language model's space};
  \node[block, right=10mm of proj] (dec) {\textbf{Decoder}\\[2pt]\scriptsize
    autoregressively generates text conditioned on visual tokens and the prompt};
  \node[endpoint, right=4mm of dec] (out) {output\\tokens};

  \draw[flow] (page) -- node[onflow, above=1mm, pos=0.5] {} (vit);
  \draw[flow] (vit) -- (proj);
  \draw[flow] (proj) -- (dec);
  \draw[flow] (dec) -- (out);

  \node[onflow, anchor=south] at ($(vit.north east)!0.5!(proj.north west) + (0,2.5mm)$)
    {visual\\embeddings};
  \node[onflow, anchor=south] at ($(proj.north east)!0.5!(dec.north west) + (0,2.5mm)$)
    {visual\\tokens};

  %% The prompt enters the decoder from below: it is the second conditioning channel,
  %% and everything task-specific in this thesis lives in it.
  \node[endpoint, below=9mm of dec] (prompt) {textual\\prompt};
  \draw[flow] (prompt) -- (dec);

  \node[
    draw=horusmuted, line width=0.5pt, rounded corners=2pt, dotted,
    fill=white, text=horusink, font=\scriptsize, align=left,
    text width=4.6cm, inner sep=4pt,
    below=9mm of vit.south west, anchor=north west,
  ] {No component of this architecture is invoice-specific: all task behaviour
    is elicited by the prompt and learned from training data.};

\end{tikzpicture}
